\section[几何 Satake 对应续论]{几何 Satake 对应续论（5 月 20 日）}

\srcpage{78}

距上次见面已有一段时间，我们先回忆上次讲的内容：几何 Satake 定理。我们把它解释为量子场论 $\cA_G$ 中线算子（line operator）构成的对称幺半范畴。线算子范畴就是 $\cA_G(S^2)$；由于身处四维，算子乘积展开赋予它对称幺半结构。具体地说，这便是球面 Hecke 范畴（spherical Hecke category）$\mathcal{S}\mathrm{ph}$，即仿射 Grassmann 簇上等变层构成的范畴。

用 Tannakian 重构，我们看到：带纤维函子的 $\mathcal{S}\mathrm{ph}$ 就是 $G^\vee$-表示范畴；这里 $G^\vee$ 是 $k$ 上的约化代数群，$k$ 是我们取层所用的系数域（或环）。这正是 $G$ 的 Langlands 对偶群。有同构 $\pi_1(G)\cong Z(G^\vee)^\vee$，即 $G^\vee$ 的中心的特征标对偶；特别地，若 $G$ 单连通，则 $G^\vee$ 伴随，反之亦然。此外，选取极大环面 $T\subset G$ 给出典范的对偶环面 $T^\vee\subset G^\vee$，它也是极大的。而且 $G^\vee$ 还带有更多数据，称为钉住（pinning）\yizhu{钉住（épinglage）指 Borel 子群、极大环面与一组根向量 $x_{\alpha_i^\vee}$ 构成的三元组；它把对偶群连同这些典范数据一起钉死。}{I（本库 SY02-36 术语表 T1276）}，它其实相当刚性：例如它决定一个 Borel 子群，还给出典范的余特征（cocharacter）
\[
2\rho^\vee\colon \G_m\hookrightarrow T^\vee.
\]
这是因为我们的纤维函子不只取值于向量空间，而是取值于 $\Z$-分次向量空间（我们把它等同于上同调）。于是 Borel 子群由正的上同调次数给出。你还得到主幂零元的一个选择！这在 Satake 定理的导出版本中变得非常相关。行列式线丛 $\cL\to\Gr_G$ 是这一结构的一部分；这与 Ginzburg 的工作有关。

你可能想要、或为之欢呼这些额外结构，其中一个哲学理由是：如果你要研究群，把表示论数据唯一地钉住是有帮助的，而钉住做的正是这件事。

Langlands 对偶的例子：

\begin{itemize}
\item 若 $T$ 是环面，其 Langlands 对偶是对偶环面 $T^\vee$。
\item $\mathrm{GL}_n$ 是 Langlands 自对偶的。也就是说，对向量丛我们看不到新的群。
\item 在 A 型中，$\mathrm{SL}_n$ 单连通但有中心，所以此时 $G^\vee$ 会不同，确实得到 $\mathrm{PGL}_n$。
\item D 型 $\mathrm{SO}_{2n}$ 是 Langlands 自对偶的。
\item B 型与 C 型互为 Langlands 对偶：$\mathrm{SO}_{2n+1}$ 与 $\mathrm{Sp}_{2n}$。这有点诡异：这两个群乍看毫无关系，而且似乎没有清楚的办法从群作用做出这种对偶。这是一个重要的启发式要点：你不能真的指望靠摆弄作用就能凑出朴素或简单的 Langlands 对偶。这个 B–C 对偶说明：Langlands 对偶往往是深刻的。
\end{itemize}

\srcpage{79}

设 $C$ 是曲线，$x\in C$。于是我们得到球面 Hecke 范畴在 $\mathrm{Shv}(\BunG(C))$ 上的作用：对象 $V\in\mathcal{S}\mathrm{ph}$ 通过对应
\[
\begin{tikzcd}
& \Hck_x \arrow[dl] \arrow[dr] & \\
\BunG & & \BunG
\end{tikzcd}
\]
作推拉作用。

对 $G=\mathrm{GL}_n$，$x$ 处纤维的 $k$ 维子空间构成的空间给出嵌入 $\mathrm{Gr}(n,k)\subset\Gr_{\mathrm{GL}_n}$。于是 $\mathrm{GL}_n$ 作用在 Grassmann 簇的上同调上。这听起来没有实际上那么令人惊讶：Grassmann 簇是齐性空间，所以 $\mathrm{GL}_n$ 作用在它上面是合理的。但这其实不是作用真正的来源！真正起作用的是 $(\mathrm{GL}_n)^\vee$，而这是第 $k$ 个基本表示。

由于 $\mathcal{S}\mathrm{ph}$ 是一个「交换环」（对称幺半范畴），我们想把它在 $\mathrm{Shv}(\BunG(C))$ 上的作用作谱分解。这需要寻找 $V\in\mathcal{S}\mathrm{ph}$ 的特征对象（eigenobject）：层 $\cF\in\mathrm{Shv}(\BunG)$ 连同同构的数据
\[
V*\cF\xrightarrow{\ \sim\ } E_x(V)\otimes \cF.
\]
这里 $E_x(V)$ 是向量空间 $V$ 上的摩天层（skyscraper sheaf），这就是我们为特征层（eigensheaf）选定的「标量」概念。把这一点稍微展开，$E_x$ 必须是函子 $\Rep(G^\vee)\to\Vect$；与张量积的相容性意味着它必须是对称幺半的。因此我们可以分类可能的「特征值」$E_x$：这样的映射等价于反方向的映射 $\mathrm{pt}\to\mathrm{pt}/G$。这样的映射在同构意义下只有一个；若把同构也记录下来，$E_x$ 就等价于一个 $G^\vee$-主丛。给定 $G^\vee$-主丛 $E_x$ 与 $V\in\mathcal{S}=\Rep_{G^\vee}$，得到的向量空间是（点上的）相伴丛 $E_x\times_G V$。这个结构有点令人困惑，因为一切都在非典范地同构。

于是让 $x$ 变动：取 $C$ 的有限点集 $S$，考虑 $\bigotimes_{x\in S}\Rep_{G^\vee}$ 在 $\mathrm{Shv}(\BunG(C))$ 上的作用。我们从头到尾的动机是：这些数据都属于一个拓扑场论，而这一作用是线算子对 $\cA_G(C)$ 的作用。TFT 结构给我们这样的想法：这些线算子的作用是局部常值的。这意味着它通过因子化同调（factorization homology）$\int_C\Rep_{G^\vee}$ 的作用分解；后者在不选取 $S$ 的情况下编码这种局部常值性。事实上，
\[
\Spec\int_C \Rep_{G^\vee}=\mathrm{Map}(C,\mathrm{pt}/G^\vee)=\LocSys_{G^\vee}(C).
\]
我们把点与点之间的路径等同起来，所以只看到 $C$ 的基本群胚——这就是我们看到局部系统的原因。

\begin{remark}
局部系统的空间，至少在 Betti 情形，极其简单：就是映射 $\pi_1(C)\to G^\vee$ 模共轭。而 $\pi_1(C)$ 有由 $A$-圈与 $B$-圈给出的简单表示，所以 $C$ 上的 $G^\vee$-局部系统就是群中 $2g$ 个元素 $A_1,\dots,A_g,B_1,\dots,B_g$ 的选择（$g$ 是 $C$ 的亏格，我们假定 $C$ 连通），满足关系
\[
\prod_i [A_i,B_i]=1.
\]
特别地，$\LocSys_{G^\vee}(C)$ 相当具体而简单：它是仿射簇对 $G^\vee$ 的商。
\end{remark}

不幸的是，de Rham 情形中的类比不成立：$\LocSys_{G^\vee}(C)$ 更复杂。

总之，我们当时在讨论：如果你相信 $\cA_G(C)$ 是拓扑场论，就会得到：$C$ 的有限子集上的算子作用通过因子化同调分解。而这并不太复杂：它就是 $\QCoh(\LocSys_{G^\vee}(C))$。这是从拓扑场论得到的主要输出之一：我们附给曲面的范畴会得到某个好东西的作用。

这些想法在许多情形下是定理，取决于你对 $\BunG(C)$ 上层的理解。这些都是深刻而困难的定理。

\begin{theorem}[Gaitsgory 消没定理]
这在 de Rham 情形成立。确切地说，$\QCoh(\LocSys^{\mathrm{dR}}_{G^\vee}(C))$ 作用在 $\cD(\BunG(C))$ 上。
\end{theorem}

这很难，是几何 Langlands 纲领中最重要的结果之一。de Rham 局部系统指带平坦联络的主 $G$-丛；而这里说的「层」指 $D$-模。

\srcpage{80}

\begin{theorem}[Nadler–Yun]
这在 Betti 情形成立；此时 $\mathrm{Shv}(\BunG(C))$ 指 $\BunG(C)$ 上幂零的层构成的范畴。
\end{theorem}

这个幂零层范畴是这样得到的：从 $\Bun_T(C)$ 或 $\Bun_L(C)$（$L$ 为 Levi 子群）之类对象上的局部常值层出发，再诱导到 $G$。想法是我们想要局部常值层，因为它们最简单；但我们需要范畴在某些换群操作下封闭，特别是抛物诱导。这与 Braverman–Gaitsgory 关于几何 Eisenstein 级数的工作有关。

\begin{theorem}[Arinkin–Gaitsgory–Kazhdan–Raskin–Rozenblyum–Varshovsky {[AGK+20b, AGK+20a, AGK+21]}]
这在有限域上的限制情形成立。
\end{theorem}

这个情形在某种意义上是 de Rham 与 Betti 的交集（「核心几何 Langlands 纲领」）：任何几何 Langlands 表述所共有的东西。因此这个受限问题在任何域上都有意义。

回到我们的故事：$\mathrm{Shv}(\BunG(C))$ 在 $\LocSys_{G^\vee}(C)$ 上作谱分解；而 $\QCoh(\LocSys_{G^\vee}(C))$ 的作用等价于 $\LocSys_{G^\vee}$ 上的一个范畴层。关于它你能问的最简单问题是：局部系统 $E$ 处的纤维是什么？这就是 $E$-Hecke 特征层：数据为 $\cF\in\mathrm{Shv}(\BunG)$，并且对每个 Hecke 算子 $V\in\Rep_{G^\vee}$ 有同构
\[
V_x*\cF\xrightarrow{\ \sim\ }(V_x)_{E_x}\otimes\cF,
\]
且这些同构以相容的方式给出。也就是说，我们在每个点 $x\in C$ 上指定了一个 $G^\vee$-主丛，并希望它们拼成一个主 $G^\vee$-丛 $E$。这些是特殊函数的类比：就像 $e^{\lambda x}$ 作为特征值 $\lambda$ 的特征函数出现，现在我们把 $\lambda$ 换成局部系统并取特征对象。于是「朴素几何 Langlands 猜想」的部分期望是：这样的对象足够多，且它们张成空间，很像指数函数张成 $L^2$；这意味着可以写下 Fourier 变换。

更确切地说，朴素猜想断言 $\cA_G(C)$ 等价于 $\QCoh(\LocSys_{G^\vee}(C))$，并且这一等价把特征层与摩天层互换。这个猜想的另一种说法是：$\QCoh(\LocSys_{G^\vee})$ 在 $\mathrm{Shv}(\BunG(C))$ 上的作用是否是秩一的自由循环作用。生成元叫做 Whittaker 层 $\cW$，它扮演单位元处摩天层的角色。无论怎么表述，这个猜想都说：关于 $\cA_G(C)$ 你想知道的一切，都可由谱分解理解。

好吧，朴素几何 Langlands 猜想是错的。这并不算特别特别深刻；在它被证伪之前，人们就知道它很可能是错的。但当我们把猜想修正成真正正确的命题时，它是一盏指路明灯。而且朴素形式错得并不算特别特别离谱。

这与 $\LocSys_{G^\vee}$ 是奇异的有关。修正的一种办法是只限制到光滑轨迹，即不可约局部系统。对 $G=\mathrm{GL}_n$，这意味着只考虑不可约平坦联络；一般地，这意味着扔掉那些结构群可约化到抛物子群的联络（TODO：我想我没弄错，但不太确定。）不可约局部系统的轨迹 $\LocSys^{\mathrm{irr}}_{G^\vee}$ 是光滑叠，而且几乎是簇。事实上，若 $G$ 单连通，从而 $G^\vee$ 无中心，则自同构消失，$\LocSys^{\mathrm{irr}}_{G^\vee}$ 就是簇。因此可以尝试对这个光滑簇证明朴素几何 Langlands 猜想的类比。对 $\mathrm{SL}_2$，这归功于 Drinfeld；对 $\mathrm{GL}_n$（对某些 $n$？TODO），这归功于 Gaitsgory。由于不可约局部系统没有到抛物子群 $P^\vee\subset G^\vee$ 的约化，它们不与几何 Eisenstein 级数对话；这类局部系统有时叫做尖点局部系统（cuspidal local system）。

尖点形式（cusp form）是 Langlands 对应中最重要、最深刻的部分，即那些不通过显然构造（抛物诱导）而来的形式。对偶侧的光滑轨迹恰好是它们，这很不错；你可能会问：那为什么还要管别的？部分原因是它们最深刻，所以极难把握：你无法从更小的 $G$ 归纳地得到它们。于是「无聊的部分」在技术上反而最难：在算术 Langlands 纲领中，所有分析都在这里进入。但只有使用整个空间才能做出一致的预测，所以两部分都需要。

除朴素几何 Langlands 猜想之外，还可以表述电磁对偶的朴素版本（Montonen–Olive 对偶或 S-对偶）：整个理论 $\cA_G$ 等价于 $\cB^n_{G^\vee}$，后者是真正的谱侧 $\cB_{G^\vee}$ 的某种朴素近似。这个朴素 B 侧理论就是我们所说的叠 $\mathrm{pt}/G^\vee$ 上的四维 B-模型。这意味着：对闭三维流形 $Y$，$\cB^n_{G^\vee}(Y)=\cO(\LocSys_{G^\vee}(Y))$；对闭曲面 $\Sigma$，$\cB^n_{G^\vee}(\Sigma)=\QCoh(\LocSys_{G^\vee}(\Sigma))$；对闭一维流形 $N$，也许还有 $\cB^n_{G^\vee}(N)=\mathrm{ShvCat}(\LocSys_{G^\vee}(N))$！再说一次，这并不完全正确，但它是真值很好的首次近似。

\srcpage{81}

（好吧，我们本来还可以做更傻的事：用局部算子（local operator）代替线算子来猜 $\cB_{G^\vee}$ 是什么。但 $\cA_G$ 中局部算子的代数是平凡的，所以那里没什么可说。）而我们的近似只用了线算子，却出乎意料地接近真相。我们也许能从面算子（surface operator）学到更多，不过已知的并不多。过渡到导出情形也能提供洞见，帮助修正朴素对偶：你不能只用 $\Rep_{G^\vee}$ 的导出范畴，而必须用更有趣的东西（TODO：想必我们最终用来替换 $\Rep_{G^\vee}$ 的东西带有一个 $t$-结构，其心是 $\Rep_{G^\vee}$，但我应当确认这确实成立）。

另一个可以问的问题是：对应于平凡 $G^\vee$-局部系统 $E$ 的 $\BunG(C)$ 上「那个」自守层是什么？这原来是 $\LocSys_{G^\vee}(C)$ 上最奇异、最像叠的点；平凡局部系统由平凡 $T^\vee$-局部系统诱导而来，或由任意 Levi 子群 $L^\vee\subset G^\vee$ 的平凡 $L^\vee$-局部系统诱导而来。于是你得到 A 侧所需层的一个公式，事实上对每个这样的子群各得一个公式。问题是：它们彼此不一致。

你可以试 $\BunG(C)$ 上最显然的层，即常值层 $k$。然后问：它是 Hecke 特征层吗？这个问题比你预想的更微妙：它取决于你究竟怎么定义。你可以把 $V*k$ 定义为 $k$ 的拉回限制在某个 $G(\cO)$-轨道上再前推，这给出 $\cH_V\in\mathrm{Shv}(\Gr_G)$ 的上同调与 $k$ 的张量积。这令人担忧：上同调并不集中在次数零。所以它不算 Hecke 特征层；但如果你允许特征值是分次向量空间，那还可以。这类对象叫做 Arthur 特征层：它们有非平凡的 Arthur 参数，意思是上同调分次（或 Lefschetz $\mathrm{SL}_2$-作用）是有趣的。所以这不算太糟，但如果你要找的是平凡局部系统处的摩天层，它就是错的对象。这正是我们对 $\cB_{G^\vee}$ 的朴素近似所缺少的那类东西。

例如，$\cA_G$ 理论有有趣的线算子，由 't Hooft 线算子给出；它们应当与 $\cB_{G^\vee}$ 中的 Wilson 线算子（更容易定义的那些）互换。这些 't Hooft 算子测量一个点周围的奇异性（例如单极子）。在某种意义下，从 Langlands 的视角看，你可能想用 A 侧来造 B 侧；但物理的视角说，从一开始就有 A 与 B 两个 twist，线算子相同，只是在两侧表现不同（这是 S-对偶猜想的一部分）。在 Riemann 面的层面，这一切只是说：$\BunG(C)$ 上的层与 $\LocSys_{G^\vee}(C)$ 上的层之间的等价，应当与 $\mathcal{S}\mathrm{ph}\cong\Rep_{G^\vee}$ 的几何 Satake 等同相容；而这两个对称幺半范畴分别作用在 $\mathrm{Shv}(\BunG(C))$ 与 $\QCoh(\LocSys_{G^\vee})$ 上。

接下来我们讨论去范畴化与取迹，这将使我们把它与经典 Satake 定理联系起来。计划是下周四下午 1:30 见面，细节见 Slack。之后我们可以讲导出版本。也许再往后还能讲到面算子（局部几何 Langlands）与边界条件。
